\section{Constructive scope and reading conventions}\label{sec:scope}

\subsection{The base theory}

We use \(\ipra\) in the standard first-order sense: the only individual sort is
\(\N\); every primitive-recursive operation may be named by a function symbol;
its defining equations are available; reasoning is intuitionistic; and
induction is restricted to quantifier-free formulas.  Finite sequences,
finite trees, signed integers, rationals, and tuples are therefore represented
by natural-number codes.  Bounded search and bounded sums are primitive
recursive.  A theorem about a coded finite object is internal when its proof is
an explicit derivation from those defining equations and quantifier-free
induction.

The restriction to quantifier-free induction is not an obstacle to the finite
probability identities in the article.  Each relevant invariant can be
maintained by an equation or decidable inequality.  For example, the statement
that the first \(n\) terms of two finite sums agree is an equation between two
primitive-recursive terms.  Its successor step is one use of the defining
recurrence for bounded summation.

\begin{remark}[No hidden classical quotient]
A rational number is not treated as an abstract quotient field.  It is a code
with a positive denominator, and equalities are proved by
cross-multiplication.  Likewise, ``almost surely'' on a finite weighted law is
an equation saying that the total weight of the exceptional outcomes is zero.
Both are decidable arithmetical relations.
\end{remark}

\subsection{Three meanings of a theorem}

The article mixes finite algebra, countable series, and general integration.
We use the following disciplined translations.

\begin{definition}[Exact finite theorem]
An exact finite theorem quantifies over natural codes for a finite law and
finite observables.  After denominator clearing, its conclusion is a
quantifier-free equation or inequality.  The proof uses only primitive
recursion and quantifier-free induction.
\end{definition}

\begin{definition}[Effective analytic theorem schema]
An effective analytic theorem is a uniform family of finite theorems indexed by
an output precision \(p\).  The input contains primitive-recursive approximants
and explicit moduli.  From \(p\), the proof computes a finite stage \(N(p)\),
applies exact finite reasoning there, and proves an error at most \(2^{-p}\).
For each fixed presentation, the entire checker is primitive recursive.
\end{definition}

\begin{definition}[Constructive replacement]
When the article states a classical existence theorem whose data do not contain
an algorithm or modulus, we replace it by a witnessed version.  We prove the
forward implications from the witness and record that the non-effective
converse has not been internalized.
\end{definition}

\subsection{What ``prove as many statements as possible'' means}

The theorem ledger accompanying this paper separates the article's content
into mathematical claims, examples, and historical prose.  We prove every
finite algebraic claim, including the standard inequalities and product-law
formulas.  We also prove effective versions of the main countable and
continuous claims.  Some article sentences are definitions rather than
propositions; in those cases the task is to show that the definition has a
primitive-recursive realization and interacts correctly with the finite
kernel.

We do not claim that every classical random variable or measure has an
\(\ipra\)-usable name.  A bare statement such as ``let \(X\) be an integrable
random variable'' suppresses both a representation and a convergence rate.
Our version instead says: let \(X_p\) be a primitive-recursive sequence of
finite rational simple laws with a certified \(L^1\) Cauchy bound.  This is
strictly more informative and is exactly what a proof checker can consume.

\subsection{Constructive logic in the proofs}

All finite predicates occurring below are decidable: equality and order of
coded integers and rationals, membership in a finite event, and comparison of a
bounded sum with a rational threshold.  Case distinctions on such predicates
are intuitionistically valid.  We never infer an arbitrary proposition \(A\)
from \(\neg\neg A\), and we never use an unbounded least counterexample.

Several classical-looking arguments become simpler when written pointwise.
Markov's inequality is the pointwise inequality
\(a\ind_{X\ge a}\le X\), summed with nonnegative weights.  The variance identity
is a polynomial identity.  Cauchy--Schwarz is a sum of squares.  The law of
total expectation is a finite regrouping of one double sum.  These proofs do
not require excluded middle beyond decidable comparisons already supplied by
the codes.

\subsection{Trust boundary}

The trusted mathematical boundary intended by this report is the existing
ProveIt calculus \code{Prf PRAAx} and Lean's kernel.  The supplied Python script
is deliberately outside that boundary: it exhaustively and randomly checks
many small exact instances, but it is only regression evidence.  Similarly,
ambient Lean structures in the packet are specifications until their fields
are filled by actual proofs and then translated to object-language derivations.
